\section{Theory}

This section derives the core theoretical results of the model. We focus on how stablecoins affect behavior through both the cost and information channels.

\subsection{Proposition 1: Stablecoin Demand and Misalignment}

\textbf{Proposition 1.} Stablecoin demand $q_s$ is increasing in exchange rate misalignment $\theta$.

\textit{Proof.} 

Agents choose $q_s$ to maximize expected utility. The marginal benefit of foreign currency holdings increases with $\theta$ because higher misalignment raises expected depreciation losses in the crisis state. Since stablecoins have lower transaction costs than cash, the marginal utility gain from increasing $q_s$ is increasing in $\theta$.

Formally, the first-order condition implies:
\begin{equation}
\frac{\partial U}{\partial q_s} = \frac{\partial U}{\partial c} \cdot \frac{\partial c}{\partial q_s} - C_s'(q_s) = 0.
\end{equation}

As $\theta$ increases, the marginal benefit term rises, implying a higher optimal $q_s$. \hfill $\square$

\subsection{Proposition 2: Endogenous Signal Precision}

\textbf{Proposition 2.} The precision of the public signal increases with stablecoin usage.

\textit{Proof.}

By assumption:
\begin{equation}
\sigma_s^2(q_s) = \frac{\bar{\sigma}_s^2}{1 + \rho q_s}.
\end{equation}

Thus, precision is:
\begin{equation}
\tau_s(q_s) = \frac{1}{\sigma_s^2(q_s)} = \frac{1 + \rho q_s}{\bar{\sigma}_s^2},
\end{equation}

which is strictly increasing in $q_s$. \hfill $\square$

\subsection{Proposition 3: Coordination and Crisis Probability}

\textbf{Proposition 3.} An increase in the precision of the public signal increases the probability of a coordinated attack.

\textit{Proof.}

As $\tau_s$ increases, the weight $w_s$ on the public signal increases while $w_x$ decreases. This reduces the dispersion of posterior beliefs across agents. In the global-games framework, lower dispersion increases the mass of agents whose beliefs cross the attack threshold simultaneously.

Formally, the variance of posterior beliefs decreases with $\tau_s$. This increases the slope of the aggregate attack function $R(\theta)$, making the crisis condition more sensitive to fundamentals. Therefore, the probability that the crisis threshold is exceeded increases. \hfill $\square$

\subsection{Proposition 4: Welfare Non-Monotonicity}

\textbf{Proposition 4.} Welfare is non-monotonic in exchange rate misalignment $\theta$ in the presence of stablecoins.

\textit{Proof.}

Welfare depends on two opposing forces:

\begin{enumerate}
\item \textbf{Efficiency gains:} Lower transaction costs and improved risk-sharing increase welfare
\item \textbf{Fragility costs:} Higher crisis probability reduces welfare
\end{enumerate}

At low levels of $\theta$, crisis probability is low, and efficiency gains dominate. At high levels of $\theta$, the increase in crisis probability dominates, leading to lower welfare.

Therefore, welfare is non-monotonic in $\theta$. \hfill $\square$

\subsection{Discussion}

The propositions highlight the central mechanism of the model: stablecoins improve information aggregation, but this very improvement strengthens coordination and increases fragility. This tradeoff is at the core of the paper's contribution.
